% Speaker script for talk.pdf -- 10-minute team presentation.
% Build with: pdflatex -interaction=nonstopmode script.tex  (twice)
\documentclass[11pt,a4paper]{article}

\usepackage[utf8]{inputenc}
\usepackage[T1]{fontenc}
\usepackage[margin=2.1cm,top=2.0cm,bottom=2.2cm]{geometry}
\usepackage[scaled=0.92]{helvet}
\renewcommand{\familydefault}{\sfdefault}
\usepackage{amsmath,amssymb}
\usepackage{xcolor}
\usepackage{enumitem}
\usepackage{booktabs}
\usepackage{microtype}
\usepackage{fancyhdr}

\definecolor{ink}{HTML}{1C2733}
\definecolor{accent}{HTML}{3B6FA0}
\definecolor{warm}{HTML}{C07C33}
\definecolor{mute}{HTML}{6B7785}
\definecolor{rule}{HTML}{D7DDE4}
\definecolor{wash}{HTML}{F4F6F8}
\color{ink}

\setlength{\parindent}{0pt}
\setlength{\parskip}{0.32em}
\linespread{1.00}

\pagestyle{fancy}
\fancyhf{}
\renewcommand{\headrulewidth}{0pt}
\fancyfoot[C]{\footnotesize\color{mute}\thepage}

% ---- slide heading: "SLIDE n -- title" with a time budget on the right
\newcommand{\slide}[3]{%
  \vspace{1.0em}%
  {\color{rule}\rule{\textwidth}{0.6pt}}\\[0.35em]%
  \begin{minipage}[t]{0.78\textwidth}%
    {\color{accent}\footnotesize\bfseries SLIDE #1}\\[0.15em]%
    {\large\bfseries #2}%
  \end{minipage}\hfill%
  \begin{minipage}[t]{0.20\textwidth}%
    \raggedleft{\footnotesize\color{mute}#3}%
  \end{minipage}\\[0.5em]%
}

% ---- stage direction / delivery note
\newcommand{\cue}[1]{{\color{warm}\small\textit{[#1]}}\par}

% ---- the words to say
\newenvironment{say}{\itshape}{\par}

\newcommand{\keyw}[1]{\textcolor{accent}{\textbf{#1}}}

\begin{document}

% ================================================================ header
{\color{rule}\rule{\textwidth}{0.9pt}}\\[0.9em]
{\LARGE\bfseries Speaker script}\\[0.4em]
{\large\color{mute} ``Back to entropy'' --- about 8:10 on the core path}\\[0.9em]
{\color{rule}\rule{\textwidth}{0.9pt}}

\vspace{0.6em}
\begin{tabular}{@{}ll@{}}
\textbf{Deck} & \texttt{slides/talk.pdf} --- 18 slides $+$ 5 backup slides \\
\textbf{Version} & Final paper-aligned presentation and discussion guide \\
\textbf{Target} & \textbf{under 10 minutes} of talking, then Q\&A on the remaining time \\
\textbf{Goal} & The team understands \emph{how the method works} and \emph{what we found} \\
\end{tabular}

\vspace{1.2em}
\textbf{How to use this document.} Italic text is what to say --- written to be spoken, not read
aloud verbatim. Orange brackets are delivery cues. Times are cumulative: if you are past the
stated time when you leave a slide, cut the next slide to its first sentence. The
\emph{must-land} points are marked \keyw{LAND THIS}; everything else is expendable under time
pressure.

\vspace{0.4em}
\textbf{Length, measured not guessed.} The spoken text is 1{,}205 words on the core path and
1{,}508 words in full, which at 150\,wpm is \textbf{8:00} and 10:01 respectively. The cumulative
times on each slide are the core path, which \emph{omits slides 6, 7, 10, 14, and 15} --- the pipeline
diagram, the contamination audit, and the math-domain bonus result. Those three are marked
\texttt{optional} with the time they cost if you keep them; they are the detail/bonus slides that
are the first thing to cut, not the must-land points. Re-check after any edit with
\texttt{python tools/script\_timing.py}, which checks that the core path stays under ten minutes.

\vspace{0.4em}
\setlength{\fboxsep}{9pt}
\colorbox{wash}{\parbox{\dimexpr\textwidth-2\fboxsep\relax}{%
\textbf{The four things that must land:}
\begin{enumerate}[leftmargin=1.4em,itemsep=0.15em,topsep=0.35em]
  \item Return to predictive entropy, made usable through byte normalization, fresh data,
  controlled adaptation, and contamination filtering.
  \item \textbf{Adapted BPB selects the model}; adaptation moves candidates by as many as seven
  rank positions, while relative reduction explains corpus calibration.
  \item BPB matches sentence continuation ($r=-0.975$, $\rho=0.982$ vs.\ HellaSwag), while knowledge and
  arithmetic benchmarks add complementary capability views.
  \item \textbf{We answered the circularity objection, at its real strength} (slide 17). On an
  in-domain criterion that is \emph{not} a benchmark, adapted BPB is the only one of seven
  selectors clearing chance among comparable-size candidates --- and we say plainly that the lead
  over the benchmarks is unresolved and that model size is not beaten. Conceding those two is what
  makes the first one land.
\end{enumerate}

\vspace{0.4em}
\textbf{Talk shape.} One ruler, one controlled intervention, an independent criterion, one decision
rule. \emph{The domain-stability result (slide 14) moved to the optional set to make room; it is
still in the deck and still worth using if the room wants robustness.}}}

% ================================================================ 1
\slide{1}{Title --- ``Back to entropy''}{0:00--0:13}

\cue{Do not read the title. Land the one-sentence framing and move on.}

\begin{say}
Quick version of what I have been building. The thesis: entropy, measured properly, gives us a direct
and repeatable way to choose a base model for a target domain. Under ten minutes, then questions.
\end{say}

% ================================================================ 2
\slide{2}{Return to the language-modeling objective}{0:13--0:52}

\cue{\keyw{LAND THIS} --- point 1. If you only get one slide across, it is this one.}

\begin{say}
Language models were originally evaluated by predictive entropy. Then evaluation moved toward fixed
task suites: useful, but each one measures only a slice of base-model quality, and repeated
optimization makes benchmark tuning and contamination difficult to rule out.

So we go full circle, but not back to raw perplexity. Perplexity is per token, and a token is a
different amount of text in every model. Comparing it across tokenizers is like comparing euros and
yen without an exchange rate. We make entropy useful again with four pieces: byte normalization,
fresh target-domain text, controlled adaptation, and contamination filtering.
\end{say}

\cue{The euros-to-yen line is the one people repeat afterwards. Deliver it slowly.}

% ================================================================ 3
\slide{3}{Why not just read the scorecard?}{0:52--1:50}

\cue{\keyw{LAND THIS} --- this is the reason the method exists. Give the table three seconds of
silence before you say the numbers. Do not rush to the punchline.}

\begin{say}
The obvious question first: every model ships with a scorecard, so why measure anything ourselves?

Look at these two. On our ruler they are three per cent apart; on commonsense, two points apart.
Essentially the same model, twice. On grade-school maths, one scores six per cent and the other
sixty-one.

That gap is not telling you one model is better --- the other two rows say they are tied. It tells
you what went into the two training mixes. Both are legitimate choices; the score cannot tell you
which one you are buying.

And note what I am not claiming. I am not saying anyone games their benchmarks --- I cannot see
anyone's training data. The point is structural: a years-old public benchmark is a target whether or
not anybody aims at it, and text written after a model ships cannot have been studied for.
\end{say}

\slide{4}{Measure the model we will actually use}{1:50--2:21}

\cue{\keyw{LAND THIS} --- point 2.}

\begin{say}
A static benchmark describes the model as released. A product uses the model after it has seen some
domain data. So we score both states on the same held-out corpus: zero-shot BPB, then adapted BPB.

The adapted level is the model-selection number because it measures the system we would actually
deploy. The reduction between the two states has a different job: it tells us how much calibration
the base distribution needed. One trajectory, two clearly defined outputs.
\end{say}

\cue{Point to the right-hand box: that adapted model is the object being selected.}

% ================================================================ 4
\slide{5}{Two measurements, one decision}{2:21--3:04}

\cue{Do not derive the formula. Point at it, say what the denominator is, move on.}

\begin{say}
Two measurements, one decision.

Move one: change the ruler. Normalize the loss by bytes instead of tokens --- bits per byte. A UTF-8
byte is the same unit for every model on earth, so the numbers become comparable. BPB is not my
invention --- The Pile used it, Paloma uses it. It is just badly under-used, and we make it primary
rather than a footnote.

Move two: measure domain fit. For each model, on one identical corpus: zero-shot BPB, a controlled
LoRA adaptation, then adapted BPB. We rank by the adapted endpoint and keep the relative reduction
as a calibration diagnostic. That gives every number one simple job.
\end{say}

% ================================================================ 5
\slide{6}{The pipeline}{optional +0:37}

\cue{Walk the boxes left to right with your hand. Forty seconds, no more.}

\begin{say}
Here is the whole thing. Raw corpus in; contamination audit; domain adaptation; BPB out. The
ordering is strict by design: nothing reaches adaptation or scoring until it clears the
contamination gate.

The corpus is a deduplicated slice of Google News from the eighth of June this year --- about a
hundred and twenty thousand articles --- split with no overlap. Eleven base models, half a billion
to thirty-five billion parameters: mostly dense Transformers, plus one sparse Mixture-of-Experts
and one Liquid model, deliberately, to see whether the metric holds across genuinely different
architectures.
\end{say}

% ================================================================ 6
\slide{7}{Phase 1 --- why we trust the numbers}{optional +1:23}

\cue{Technical slide. Give the shape of the idea, not the detail. If you are behind schedule, say
only the first and last paragraphs.}

\begin{say}
A word on the contamination audit, because it is what makes the headline result mean anything.

The honest headline first: for three of our four corpora we do not need a detector at all. Reddit,
Hacker News and the math corpus were all scraped after every model in the cohort shipped, so they are
clean by construction. Freshness is a stronger contamination defence than any screen, and it is the
one we lean on.

The news corpus is the exception --- June 2026, which overlaps part of the cohort's release window
--- so it gets the screen. Scout-then-screen: cheap detectors propose suspects, a perturbation test
adds evidence. Two scouts: Min-K++, which flags odd confidence on a passage's rarest tokens, and
CoDeC, which is ours --- feed the model the text immediately preceding a passage and watch for a
sharp jump in confidence. Then a quiz that hides the real passage among two perturbed variants.

Be careful how you say the last part: preference for the original is a risk signal, not proof of
memorization --- a model can prefer the fluent original having never seen it. We quarantine the union
across all models, eighty sequences. No finite screen rules out pre-training exposure; only exact
overlap is direct evidence.
\end{say}

% ================================================================ 7
\slide{8}{Results --- 11 models, one ruler}{3:04--3:44}

\cue{The money slide. Read the axis once, then three models. Do not narrate all eleven. The exact
table is in backup if anyone wants digits.}

\begin{say}
Results. Eleven models, one ruler. One row per model, grey dot as released, blue dot after
adaptation, further left is better; the label is the reduction.

Three readings. Gemma-4-31B ends furthest left at 0.624, the best adapted result. The
thirty-five-billion Mixture-of-Experts lands with dense models in the twelve-to-fourteen-billion
range, which gives us a useful quality measurement between its active and total parameter counts.
And the Liquid model shows the largest calibration shift, 36.4 percent.

The endpoint chooses the model. The line length explains how much its base distribution moved toward
the corpus. Keeping those roles separate makes the chart straightforward.
\end{say}

% ================================================================ 8
\slide{9}{Two numbers, two jobs}{3:44--4:13}

\cue{Keep this operational: endpoint for selection, reduction for explanation.}

\begin{say}
The chart contains two useful numbers.

Adapted BPB is the primary selection score. It measures held-out target-domain prediction after the
same kind of intervention we would use in deployment, and lower is better.

Relative reduction is the calibration diagnostic. It shows how much the released model adjusted to
the corpus's style, formatting, and context pattern. It explains the path, while the endpoint makes
the decision. We verified that interpretation token by token.
\end{say}

\cue{Pause on ``token by token'' before advancing.}

% ================================================================ 9
\slide{10}{Context explains the calibration signal}{optional +1:01}

\cue{\keyw{LAND THIS} --- this is the new result and the answer to their question. Point at the
crossing. Do not rush.}

\begin{say}
We scored every token of the held-out set under the base and adapted models, then looked at where
probability moved. Two patterns make the calibration interpretation concrete.

First, the picture. We aggregate corpus loss over independent five-hundred-and-twelve-token blocks,
so almost every cut starts mid-article without the preceding context. The plot is byte-normalized,
so tokenizer vocabulary size is no longer the simple explanation --- and the curves still cross.
Over the first eight positions Gemma averages 2.566 bits per byte against Qwen's 1.490; from position
two-fifty-six on, Gemma is at 0.717 and Qwen at 0.773. With context, Gemma shows the stronger
steady-state result. Adaptation quickly learns that context-recovery pattern.

Second, which tokens gain: full stops, commas, paragraph breaks, "the", "and", and "a". Within the
top two hundred gain-carrying tokens, formatting, punctuation, and function words account for
eighty-three percent of the mass. That is a clear signature of calibration to the corpus's surface
form.
\end{say}

\cue{Land the reporting rule: adapted level for selection, reduction for calibration.}

% ================================================================ 10
\slide{11}{Adaptation does not improve the ranking --- it rotates it}{4:13--5:02}

\cue{\keyw{LAND THIS} --- this is the central finding of the talk. Read the little table left to
right, one row at a time. Do not rush it.}

\begin{say}
Now compare both rankings --- before and after adaptation --- against the three benchmarks.

Zero-shot news BPB agrees with GSM8K at 0.84 and MMLU-Pro at 0.86. After adaptation those fall to
0.73 and 0.76, while HellaSwag goes the other way, 0.72 up to 0.98. The same three-way pattern
repeats on Reddit and Hacker News.

So adaptation did not find a better model --- it rotated the ranking, toward continuation quality
and away from the generative tasks. And HellaSwag agrees so well because it scores continuations by
log-likelihood, the same object BPB measures: convergent validity, not proof.

Concretely: zero-shot BPB picks Qwen-3.5-35B, the actual GSM8K and MMLU leader; adapted BPB picks
Gemma-4-31B, the HellaSwag leader, two to four points behind on the other two.
\end{say}

% ================================================================ 11
\slide{12}{Different tasks reveal different strengths}{5:02--5:50}

\cue{\keyw{LAND THIS} --- the centrepiece. Slow down. Let people read the table for a beat before
you talk over it.}

\begin{say}
This is the table I find most useful in the whole project.

Eleven models. Left column is their BPB order; then each benchmark's own ordering of the same eleven,
with a marker for how far each model moves.

The bottom two rows are the result. HellaSwag: seven of eleven in exactly the same position, and no
model moves more than one place. MMLU-Pro, knowledge, drops to three of eleven. GSM8K, arithmetic,
is zero of eleven --- not one model keeps its place, and Gemma-4-12B moves five places.

The columns are complementary: BPB isolates target-domain language modelling, while the benchmarks
add task-specific capability information. The trend is clear, but we do not claim task complexity
forms a one-dimensional causal ladder.
\end{say}

\cue{The 7/11, 3/11, 0/11 progression is the sentence to land before advancing.}

% ================================================================ 12
\slide{13}{The family pattern explains the shift}{5:50--6:26}

\cue{The two-model box on the right is the memorable example --- point at it.}

\begin{say}
The GSM8K shift follows a strong model-family pattern. Every Qwen is unchanged or moves up relative
to its BPB position, while both Gemmas and Llama move down. This is consistent with different
math/reasoning emphasis, although this observational comparison cannot isolate a causal training ingredient.

The box on the right is the sharpest version. Qwen-2.5-1.5B and Llama-3.2-1B are adjacent in BPB,
and differ by only two points on HellaSwag. On GSM8K they score 0.611 and 0.065: a 54.6-point gap.

Same two models, two capabilities. The preferred model follows the capability the product needs.
\end{say}

% ================================================================ 13
\slide{14}{Adaptation changes the ranking --- then it transfers}{optional +1:00}

\cue{\keyw{LAND THIS} --- adaptation changes the choice; the adapted result then transfers.}

\begin{say}
Last result: all eleven models on news, Reddit, and Hacker News under one identical adaptation budget.
First, adaptation materially changes the choice. Zero-shot versus adapted rank correlations are only
0.655, 0.609, and 0.682, with maximum moves of seven, seven, and five places. Gemma-4-31B moves from
fourth, fifth, and sixth zero-shot to first adapted on news, Reddit, and Hacker News. So the adapted
endpoint is not a rescaled zero-shot ranking.

Second, that adapted result transfers. Each panel compares two corpora; Pearson correlations are
0.992, 0.977, and 0.990, while Spearman ranges from 0.955 to 0.991. Gemma-4-31B is first on all three,
and nothing moves more than two places across domains.

Every rank change sits inside a near-tied cluster. On Hacker News, Gemma-4-12B is only eight
thousandths behind Qwen-MoE and nine behind Ministral. So adapted BPB identifies the winner and stable
tiers; a target-domain run resolves the choice inside a near-tie.
\end{say}

\cue{If asked about the MoE: three billion active parameters of thirty-five billion, and it scores
like a dense 12--14B. BPB gives the direct quality measurement.}

% ================================================================ 14
\slide{15}{Domain-matched adaptation narrows the GSM8K gap}{optional +1:06}

\cue{Cut this first under time pressure --- it is a bonus result, not one of the four must-land
points. If you keep it: this is the direct test of the open question from the GSM8K slide earlier.}

\begin{say}
If the rotation is real, it should be steerable --- so we steered it. Same protocol, but adapt on
math-domain prose instead of news: arXiv abstracts, not GSM8K questions, screened against the GSM8K
test set before any model saw it.

GSM8K agreement goes to 0.94, the best selector in the study --- fifty of fifty-five model pairs
ordered correctly --- and the top pick goes back to being the actual GSM8K leader. It demotes exactly
the Gemma models GSM8K already penalizes.

The controls are the important part. Reddit gives 0.69, Hacker News 0.79, news 0.73 --- all at or
below the 0.84 you get for free by not adapting at all. Only the topically matched corpus beats
doing nothing. And it is not free either: HellaSwag agreement drops from 0.98 to 0.82.

If someone asks about significance: at eleven models the bootstrap interval on that difference still
spans zero. The evidence is the pattern across four corpora, not any single contrast. Say that
before they do.
\end{say}

% ================================================================ 15
\slide{16}{We tried to break it. The ranking held.}{6:26--6:54}

\begin{say}
Everything so far infers what adaptation does from how the \emph{ranking} moves, so we kept the
adapters and scored the adapted models on the benchmarks directly. The ordering came back identical
in all twenty cells --- not one inverted pair, which under reshuffling happens about once in a
hundred and twenty times. And re-scored as head-to-head picks with uncertainty over models, the
metric beats chance in twenty-three of twenty-four settings.
\end{say}

% ================================================================ 16b
\slide{17}{The hardest objection --- and the criterion we built}{6:54--8:21}

\cue{\keyw{The credibility slide.} Do not rush it and do not oversell it --- the concession is what
makes the rest land. If you are behind, cut the second paragraph, never the third.}

\begin{say}
Now the objection you should be forming: everything so far is validated against the same benchmarks
whose gameability is our whole argument.

So we built a criterion that is not a benchmark: items from our own held-out split, predict a
hidden name or number. The model must \emph{generate} it and match the string --- never a
likelihood score, which would agree with BPB by construction.

We built one for all four corpora and asked each which selector predicts it, among candidates
within a factor of two in size.

Start with how little is resolvable. Across four corpora, two scoring rules and seven selectors,
exactly one interval excludes chance --- ours, on news. On all four, \emph{not one public benchmark
clears chance} under either rule; neither does the free zero-shot number, nor picking the bigger
model. At ten pairs almost nothing separates, so this is a failure to reject, not proof the suites
are useless --- but it is the scorecard failing to resolve the decision it gets used for.

So we enlarged the cohort: nine more base models, preregistered; six landed, taking in-band pairs
from eleven to forty. There adapt-and-rank beats counting parameters by twenty-five points,
interval seven to fifty --- the first resolved comparison in the study. Under the prefix rule it
spans zero, so: convention-specific.
\end{say}

% ================================================================ 16
\slide{18}{A simple decision rule}{8:21--9:50}

\begin{say}
What we learned: BPB gives one coordinate system across tokenizers, and adaptation rotates the
ranking toward whatever the corpus exercises rather than improving it. The corpus is not a setup
detail --- it is the question you are asking.

How to use it. Collect recent target-domain text. More than about a factor of two apart in size?
Take the bigger one --- we are not better than that. Close? Adapt each finalist on your text and rank
by adapted BPB: minutes on one GPU. Match the corpus to the capability that decides it.

And one correction we made to ourselves, with a reason rather than a p-value. We used to recommend
the free zero-shot number. The problem is that it measures two things at once: how well a model
knows your domain, and how far its writing conventions sit from yours. Look at LFM: zero-shot it is
last of eleven, behind a model a third its size. Adapt it and it moves ahead of that model --- and
on our own criterion it really is ahead. The free number read a style mismatch as incapacity. The
three models it under-rates most are exactly the three that gain most from adaptation. It is a
screen now.

The bottom line: one fair ruler, a screen and a selector, and a corpus you choose deliberately.
\end{say}

% ================================================================ 17
\slide{19}{Thank you / questions}{9:50--9:59}

\begin{say}
That is the tour. Happy to go deeper --- backup slides cover the benchmark scores, contamination
numbers, training protocol, and the longer-context result.
\end{say}

\vspace{1.4em}

% ================================================================ Q&A
{\color{rule}\rule{\textwidth}{0.9pt}}\\[0.7em]
{\Large\bfseries Discussion guide}\\[0.2em]
{\small\color{mute} Short answers to the most useful follow-up questions.}\\[0.5em]
{\color{rule}\rule{\textwidth}{0.9pt}}

\vspace{0.5em}

\textbf{``What is BPB in one sentence?''}\\
BPB is cross-entropy divided by UTF-8 bytes, so every tokenizer and architecture is measured in the
same unit on the same text.

\textbf{``What is new in this work?''}\\
The contribution is the complete workflow: fresh domain data, contamination screening, controlled
adaptation, adapted-BPB model selection, token-level interpretation, and three-domain validation.

\textbf{``Why rank by adapted BPB?''}\\
It measures the model after the same kind of lightweight domain intervention used in deployment. It
is therefore the closest score to the system we would actually ship.

\textbf{``What does the relative reduction tell us?''}\\
It is a calibration diagnostic. The token analysis shows that formatting, function words, and
context recovery carry most of the movement, so the reduction explains the path to the endpoint.

\textbf{``Does 8.6\% of Gemma's gain come from the first 17 block positions?''}\\
No. It is the injected document-start marker plus the next 16 tokens after each document boundary,
which may occur anywhere inside a packed block. We call this the post-boundary window; the marker
itself is masked from the final cross-domain BPB.

\textbf{``Could Gemma's vocabulary explain the large early loss?''}\\
Vocabulary size can affect nats per token, so the final plot uses bits per UTF-8 byte. The crossing
remains after byte normalization, showing that context-start sensitivity is the stronger explanation.

\textbf{``Does the 512-token adapter generalize to longer contexts?''}\\
Yes in our fixed-target check. Gemma's base BPB continues improving from 0.6465 at 512 tokens to
0.6054 at 2,048, while adapted BPB stays essentially flat at 0.547. The adapter therefore retains a
9.6 percent advantage at four times its training context. Qwen's smaller gain also remains positive.

\textbf{``How does the contamination audit work?''}\\
Min-K++ and CoDeC propose suspicious passages; the Dynamic Choice Quiz tests preference for the
original over perturbations. We quarantine the union of risk-flagged passages across models so every
candidate sees one identical screened corpus. This is conservative filtering, not proof of memorization.

\textbf{``How do the static benchmarks complement BPB?''}\\
HellaSwag confirms the sentence-continuation signal ($r=-0.975$, $\rho=0.982$). MMLU-Pro and GSM8K
add knowledge and arithmetic views; their rank agreements with BPB are 0.764 and 0.727.

\textbf{``Why does the rank table use eleven models?''}\\
The corrected analysis uses the complete cohort across BPB and all three static tasks, avoiding the
stale nine-model subset that appeared in an earlier draft.

\textbf{``How stable is the result across domains?''}\\
Pairwise Spearman correlations range from $+0.955$ to $+0.991$ across news, Reddit, and Hacker News.
Gemma-4-31B ranks first on all three, and no model moves more than two places.

\textbf{``Does adaptation actually change which model we select?''}\\
Yes. Zero-shot versus adapted rank correlations are only $0.655$, $0.609$, and $0.682$ across the
three corpora, with maximum moves of seven, seven, and five places. Gemma-4-31B moves from ranks
four, five, and six to first on all three after adaptation.

\textbf{``What is the practical selection rule?''}\\
Rank candidates by adapted BPB on recent target-domain text. Treat models within $0.01$ BPB as a
near-tied tier and use a domain-specific run to make the final choice.

\textbf{``What does the MoE result mean?''}\\
The 35B MoE has about 3B active parameters and scores alongside dense 12--14B models. BPB gives a
direct quality measurement when neither active nor total parameter count is sufficient.

\textbf{``What does this cost to run?''}\\
The shared-budget comparison uses lightweight LoRA adaptation rather than full fine-tuning. It fits
the existing Slurm allocation and produces one reusable score from one domain corpus.

\end{document}
